\documentclass[conference]{IEEEtran}
\IEEEoverridecommandlockouts

% Packages
\usepackage{cite}
\usepackage{amsmath,amssymb,amsfonts}
\usepackage{algorithmic}
\usepackage{graphicx}
\usepackage{textcomp}
\usepackage{xcolor}
\usepackage{hyperref}
\usepackage{booktabs}    % For professional, clean tables
\usepackage{microtype}   % For perfect typographical alignment and spacing
\usepackage{cleveref}    % For smart, automated cross-referencing
\usepackage[spanish]{babel}
\usepackage[utf8]{inputenc}

\def\BibTeX{{\rm B\kern-.05em{\sc i\kern-.025em b}\kern-.08em
    T\kern-.1667em\lower.7ex\hbox{E}\kern-.125emX}}

\begin{document}

\title{Perfilado Automático de Complejidad de Hardware para Modelos YOLO en Despliegues de Borde\\
\thanks{Con el apoyo de eCaptureDtech.}
}

\author{\IEEEauthorblockN{William Steve Rodriguez Villamizar (wisrovi rodriguez)}
\IEEEauthorblockA{\textit{Líder de IA y Arquitecto de Soluciones} \\
\textit{eCaptureDtech}\\
Badajoz, Extremadura, España \\
wisrovi.rodriguez@gmail.com}
}

\maketitle

\begin{abstract}
El despliegue de modelos de detección de objetos en dispositivos de borde (edge) exige un cumplimiento estricto de las restricciones de hardware, específicamente en lo que respecta a la latencia, el consumo de memoria de video (VRAM) y las operaciones de punto flotante (GFLOPs). Presentamos el \textit{ModelComplexityProfiler}, un módulo de perfilado automático integrado directamente en un pipeline MLOps distribuido. Este módulo se ejecuta inmediatamente después del entrenamiento, evaluando el costo computacional arquitectónico y empírico del modelo YOLO en múltiples resoluciones de imagen. Al combinar métricas teóricas de complejidad (GFLOPs, recuento de parámetros) con mediciones empíricas (latencia en el mundo real y asignación pico de VRAM a través de \texttt{ptflops} y \texttt{pynvml}), el perfilador proporciona un plan de acción determinista que guía las decisiones posteriores de poda (pruning) y cuantización. Evaluamos el perfilador en arquitecturas YOLOv8 micro, nano y small, demostrando cómo el perfilado automático in situ acelera el ciclo de despliegue desde los servidores de entrenamiento hasta los dispositivos de borde.
\end{abstract}

\begin{IEEEkeywords}
IA de Borde, Complejidad de Hardware, YOLO, MLOps, GFLOPs, Perfilado de Latencia
\end{IEEEkeywords}

\section{Introducción}
El despliegue ubicuo de sistemas de Visión por Computadora en entornos de borde—como drones autónomos, cámaras de vigilancia inteligentes y dispositivos móviles—ha exacerbado la tensión entre la precisión algorítmica y las capacidades del hardware. Aunque los detectores de objetos modernos como YOLO logran un estado del arte en Precisión Media (mAP), su sobrecarga computacional a menudo viola los estrictos límites térmicos, de batería y de memoria del hardware de borde.

Tradicionalmente, optimizar un modelo para el borde es una fase manual y aislada, desacoplada del proceso de entrenamiento. Los ingenieros comprimen iterativamente el modelo mediante poda o cuantización y posteriormente lo perfilan en el hardware de destino. Este flujo de trabajo inconexo introduce una latencia significativa en el ciclo de vida MLOps. Para mitigar esto, proponemos una arquitectura de perfilado automático que evalúa la complejidad del hardware \textit{in situ}, inmediatamente después de que el modelo completa su fase de entrenamiento principal.

\section{Trabajos Relacionados}
La búsqueda y perfilado de arquitecturas neuronales conscientes del hardware han ganado una tracción significativa. Herramientas como \texttt{ptflops} y el \texttt{nvml} de NVIDIA proporcionan interfaces estándar para evaluar GFLOPs y memoria \cite{dollar2021fast}. Sin embargo, estas métricas generalmente se analizan de forma aislada. Integrarlas en un pipeline MLOps distribuido y automatizado—donde las métricas dictan instantáneamente la viabilidad del despliegue—representa una evolución crítica en la automatización de la IA de Borde.

\section{Arquitectura Propuesta / Metodología}
El \textit{ModelComplexityProfiler} es un módulo determinista dentro del marco más amplio de post-entrenamiento WPipe. Se invoca automáticamente tras el entrenamiento exitoso de un modelo YOLO.

\subsection{Métricas Teóricas}
Usando \texttt{ptflops}, el módulo analiza estáticamente el grafo computacional de PyTorch para extraer el recuento total de parámetros y GFLOPs. Esto proporciona una línea base independiente del hardware sobre la capacidad del modelo.

\subsection{Perfilado Empírico}
Las métricas teóricas a menudo no logran capturar los cuellos de botella de ancho de banda de memoria o los matices de optimización de la arquitectura CUDA subyacente. Por lo tanto, el módulo realiza un perfilado empírico:
\begin{itemize}
    \item \textbf{Latencia:} El modelo infiere un lote estandarizado de imágenes de validación sintéticas y reales, con tiempos registrados mediante temporizadores de alta precisión.
    \item \textbf{Asignación de VRAM:} Usando \texttt{pynvml}, el módulo rastrea la huella de memoria pico de la GPU durante el paso hacia adelante (forward pass).
\end{itemize}

\subsection{Escalado de Resolución}
Debido a que los dispositivos de borde escalan dinámicamente las resoluciones de entrada para conservar energía, el perfilador itera a través de resoluciones estándar (ej., 320, 640, 1024), generando una matriz de costo de hardware que cruza el tamaño de la imagen con la latencia y la VRAM.

\section{Configuración Experimental y Detalles de Implementación}
El perfilador opera dentro de un contenedor Docker efímero en un nodo trabajador RTX 3060 (12\,GB). Evaluamos el módulo en tres arquitecturas: YOLOv8n, YOLOv8s y YOLOv8m. El pipeline genera automáticamente un reporte JSON que contiene la matriz de complejidad, el cual luego es procesado por el LLM del pipeline.

\section{Resultados y Discusión}
El perfilado automático reveló un escalado no lineal en el consumo de VRAM a medida que aumentaban las resoluciones de imagen. Por ejemplo, mientras que los GFLOPs teóricos escalaban cuadráticamente con la resolución, la latencia empírica exhibió incrementos en función escalón debido a fallos de caché y saturación del ancho de banda de memoria en la GPU. Al exponer automáticamente estos cuellos de botella, el pipeline marcó de inmediato el modelo YOLOv8m como no viable para un dispositivo de borde con 2GB de VRAM a una resolución de 640x640, provocando un descenso (downgrade) automático a YOLOv8s.

\section{Impacto Más Amplio / Declaración Ética}
Automatizar el perfilado de hardware contribuye a la democratización de la IA de Borde, permitiendo a los investigadores optimizar modelos sin requerir acceso físico a un conjunto de dispositivos de borde. Además, optimizar modelos para una menor complejidad computacional reduce directamente la huella de carbono de las redes de sensores desplegadas.

\section{Declaración de Disponibilidad de Datos y Código}
Este ecosistema opera bajo un Modelo de Licencia Dual (PolyForm Noncommercial / AGPLv3). 
El código está disponible en \url{https://github.com/wisrovi/}.

\section{Conclusión y Trabajos Futuros}
El \textit{ModelComplexityProfiler} acorta la brecha entre el entrenamiento del lado del servidor y el despliegue en el borde. Al automatizar la extracción de métricas de hardware teóricas y empíricas, aceleramos la transición del diseño del modelo a la aplicación en el mundo real. Los trabajos futuros integrarán este perfilador en un bucle de retroalimentación que ajuste automáticamente los hiperparámetros de la arquitectura de red para satisfacer un presupuesto de hardware predefinido.

\section*{Agradecimientos}
El autor desea agradecer a eCaptureDtech por proporcionar la infraestructura y el entorno necesarios para esta investigación.

\bibliographystyle{IEEEtran}
\bibliography{references}

\end{document}
